\section{Related Work}
\label{sec:related}
\paragraph*{Efficient object detection}

YOLO evolved from single-network detection \cite{redmon2016yolo} through
multi-scale prediction and feature/training refinements
\cite{redmon2018yolov3,bochkovskiy2020yolov4}, anchor-free heads and assignment
\cite{ge2021yolox}, deployment-oriented structures \cite{li2022yolov6},
training and aggregation improvements \cite{wang2023yolov7,wang2024yolov9},
and NMS-free prediction \cite{wang2024yolov10}. These developments motivate
YOLO11 as a representative detector; GPU benchmarking does not establish
complete placement on DLA. YOLO11's detection towers and distributional
regression are retained \cite{jocher2024yolo11,li2020gfl}. Its surrounding
structure draws on feature partitioning and pyramid aggregation
\cite{wang2020cspnet,lin2017fpn,liu2018panet}. Unlike RepVGG-style deployment
re-parameterization \cite{ding2021repvgg,li2022yolov6}, the proposed adaptation changes
the incompatible attention computation while preserving this surrounding
structure. The objective is accelerator-compatible deployment rather than
optimization of the AP/FLOPs trade-off.

\paragraph*{Hardware-aware and edge network design}

YOLO-ReT and YOLOBench address edge-GPU detection and controlled evaluation
across embedded hardware \cite{ganesh2022yoloret,lazarevich2023yolobench}.
ActNAS examines hardware-sensitive activation choices, including compiler
effects \cite{sah2025actnas}. MnasNet and ProxylessNAS incorporate target
latency into architecture selection \cite{tan2019mnasnet,cai2019proxyless};
Once-for-All supports deployment specialization \cite{cai2020ofa}; MCUNet
co-designs networks and runtime \cite{lin2020mcunet}. Here strict compilation
is a feasibility constraint for a manually specified adaptation. C2DLA uses
projection and normalization related to external attention
\cite{guo2023external}, but only mixes channels at each position and adds
local depthwise context. It does not reproduce full external attention or
claim representational superiority.

\paragraph*{Heterogeneous accelerator deployment for embedded perception}
Archet et al.\ \cite{archet2023energy} study CNN design and inference options
on Jetson AGX Orin, showing that its CPU, GPU, and DLA offer different latency
and energy trade-offs. Hern\'andez et al.\ \cite{hernandez2024optimizing}
evaluate quantization on AGX Orin and report that DLA layer incompatibilities
degrade inference performance and efficiency in their tested models.
These studies make network compatibility part of the practical
accelerator-selection problem.

Tayal and Simmhan \cite{tayal2024multiinstance} examine multi-instance inference
across Orin's CUDA cores, Tensor Cores, and DLAs, investigating resource
congestion under concurrent workloads. Chakraborty et al.\
\cite{chakraborty2025profiling} profile concurrent vision inference on Jetson,
linking throughput degradation to resource contention and CPU--GPU scheduling;
aggregate GPU utilization alone does not describe component-level efficiency.
These findings motivate explicit workload placement, without establishing
concurrent-task benefits for the adapted detector.

Compiler lowering \cite{chen2018tvm} and TensorRT's optional GPU fallback
\cite{nvidia_dla_working} make compatibility an architectural and deployment
property, not merely an exporter setting. The proposed architectural adaptation
and static prediction boundary support fallback-disabled compilation and
optimized-engine placement verification, evaluated separately from deployed
accuracy, latency, host-side overhead, and module energy.
